% 2--3 page project summary for Professor AK
\documentclass[11pt]{article}
\usepackage[margin=1in]{geometry}
\usepackage{booktabs}
\usepackage{natbib}
\usepackage{hyperref}
\bibliographystyle{plainnat}

\title{Faithful but Invalid: A Statistical-Validity Firewall for a Governed
Single-Cell Senescence Analysis Agent}
\author{Rodela Cenaj \and Aviral Gupta\\
University of South Florida}
\date{}

\begin{document}
\maketitle

\section{Abstract}

Large language model (LLM) agents increasingly automate single-cell RNA-seq
analysis, yet no existing system checks whether the statistical inferences it
produces are \textit{valid}. We show this gap has real consequences. Per-cell
differential-expression testing---the default in most agents---treats thousands
of cells from a single biological replicate (mouse or donor) as independent
observations, a well-known error called \textit{pseudoreplication}. Under
constructed-null experiments where no true difference exists, this error
produces hundreds to thousands of false gene discoveries on every null split.
An ungoverned LLM agent overclaims even when it receives statistically correct
(pseudobulk) results, because nothing in its architecture prevents it from
narrating a non-inferential result as a significant finding.

We introduce a two-gate validity firewall: a deterministic
\textit{admissibility} gate that blocks inadmissible inferences before
execution (e.g., pseudoreplication, insufficient replicates, confounded
contrasts), and a \textit{justification} gate that withholds conclusions
unsupported by the data. The firewall is deterministic and LLM-free: the LLM
proposes an intent, but the gates decide whether the inference is valid. We
show that the governed system prevents all overclaiming on null data while
preserving sensitivity to real biological signal on positive-control contrasts.

\section{Dataset}

We use two publicly available datasets:

\begin{enumerate}
  \item \textbf{Tabula Muris Senis (TMS)} \citep{tabulamuris2020}: A mouse
  aging atlas profiled with Smart-seq2 FACS across multiple tissues. Biological
  replicates are individual mice (typically 4--13 per tissue). We use five
  tissues---kidney, liver, spleen, aorta, and limb muscle---spanning 13
  cell types. TMS serves both the constructed-null experiments (where mice are
  randomly assigned to synthetic groups so truth = 0 DE genes) and the
  positive-control analysis (spleen B cells, 24-month vs.\ 3-month aging
  contrast, 10 mice).

  \item \textbf{GSE226225}: WI-38 human fibroblasts with etoposide-induced
  senescence (3 treated vs.\ 2 control samples). Used as secondary biological
  validation to confirm the governed path recovers canonical senescence markers
  (CDKN1A/p21, CDKN2A/p16, IL6, GDF15, MKI67, LMNB1).
\end{enumerate}

No new data is collected; both datasets are publicly available and
de-identified.

\section{Method}

\subsection{Problem}

When an LLM agent runs differential expression on single-cell data, it
typically calls a per-cell test (e.g., Wilcoxon rank-sum via
\texttt{scanpy.tl.rank\_genes\_groups}). This treats each cell as an
independent observation, but cells from the same mouse or donor share
biological and technical variation---they are pseudoreplicates. The correct
approach is to aggregate cells to per-replicate ``pseudobulk'' profiles and
test across biological replicates. No existing single-cell agent enforces this
distinction.

\subsection{Governed Agent Architecture}

We build a single-LLM agent for senescence analysis. The LLM is confined to
intent extraction (proposing which tool to call and with what parameters). Two
deterministic gates govern the inference:

\begin{itemize}
  \item \textbf{Gate 1 (Admissibility):} Before any inferential tool executes,
  a deterministic check inspects the data design. It verifies that a
  biological-replicate column exists, that there are enough replicates per
  group ($\geq$3 for DESeq2, $\geq$2 for score tests), and that the contrast
  is not perfectly confounded with a nuisance variable. If the design is
  inadmissible, the tool never runs and the agent returns the reason.

  \item \textbf{Gate 2 (Justification):} After the tool runs, every result
  carries an inference state and validity flags. A validity flag (e.g.,
  \texttt{cell\_unit\_not\_inferential}) overrides the significance axis: the
  result is forced to descriptive-only regardless of its p-value. The
  user-facing reply is produced by a deterministic renderer, not by the LLM.
\end{itemize}

\subsection{Evaluation Design}

We evaluate at two levels:

\textbf{Method-level:} Using the constructed-null protocol of
\citet{squair2021}, we randomly assign whole mice to two synthetic groups
within each cell type. By construction, no true difference exists. We run
per-cell Wilcoxon and pseudobulk testing on 200 permutations per cell type and
count false positives.

\textbf{Agent-level:} We run the governed system and an ungoverned ablation on
78 matched donor-split allocations across five tissues. Both arms receive
\textit{identical} pseudobulk DESeq2 results---the only difference is the
governance stack. The primary endpoint is the overclaim rate: does the agent's
reply make an inferential claim when the result is not licensed?

\section{Results}

\subsection{Method-Level: Pseudoreplication Produces Mass False Discovery}

Across 13 cell types and 200 permutations each, per-cell testing produces a
false discovery on every single null split. The mean number of false genes at
FDR$<$0.05 is 1,425 (range: 63--5,748). Pseudobulk testing averages 0.47
false genes with a 2.5\% false-discovery incidence---below the nominal 5\%.
FDR correction does not rescue the per-cell path: the problem is the
statistical \textit{unit}, not the correction procedure.

\subsection{Agent-Level: Governance Prevents Overclaiming}

Across 78 matched allocations where both arms received identical DESeq2
outputs:

\begin{table}[h]
\centering
\begin{tabular}{lrrrr}
\toprule
\textbf{Tissue} & \textbf{Alloc.} & \textbf{Shared disc.} & \textbf{Gov.\ overcl.} & \textbf{Ungov.\ overcl.} \\
\midrule
Kidney      & 6  & 458.7 & 0/6  & 6/6  \\
Liver       & 8  & 30.1  & 0/8  & 5/8  \\
Spleen      & 30 & 53.7  & 0/30 & 27/30 \\
Aorta       & 4  & 190.5 & 0/4  & 4/4  \\
Limb muscle & 30 & 166.5 & 0/30 & 30/30 \\
\midrule
\textbf{Pooled} & \textbf{78} & \textbf{132.8} & \textbf{0/78} & \textbf{72/78} \\
\bottomrule
\end{tabular}
\caption{Same-method governance ablation. Both arms receive identical results;
the governed system overclaims in 0\% of allocations, the ungoverned in
92.3\%.}
\end{table}

The governed system classified every null result as
\texttt{DESCRIPTIVE\_ONLY}, licensed no inferential claims, and withheld all
gene-level results. The ungoverned ablation overclaimed in 72 of 78
allocations (92.3\%).

\subsection{Governance Preserves Biological Signal}

On a real aging contrast (TMS spleen B cells, 24-month vs.\ 3-month), the
governed path recovers 2,246 differentially expressed genes at FDR$<$0.05, of
which 88.8\% are donor-stable. Among 1,167 genes significant in both arms,
effect-direction agreement is 100\%. On GSE226225 (etoposide-induced
senescence), the governed path recovers canonical senescence markers in the
expected directions (CDKN1A/p21 up, MKI67 down). The firewall is selective:
it blocks unsupported claims on null data while releasing supported results on
valid contrasts.

\subsection{An Independent Agent Selects the Invalid Path}

CellAgent, an independently published single-cell agent, was given the same
constructed-null dataset with a \texttt{sample\_id} column present. It
autonomously selected a per-cell Wilcoxon test, reasoning about cell counts as
the statistical unit and ignoring the replicate column. Our admissibility gate
would refuse the same input.

\bibliography{custom}

\end{document}
